\documentclass[../../main.tex]{subfiles}
\begin{document}

\section{Funksjonalanalyse}
\label{sec:functional_analysis}

\subsection{Funksjonsrom}
Rommet av funksjoner som tilfredsstiller randbetingelsene.
\begin{equation}\label{eq:function_space}
V = \{ v \in C^2(\Omega) \, | \, v(a) = \alpha, v(b) = \beta \}
\end{equation}

Dette rommet representerer alle mulige funksjoner som kan være løsninger på problemet vårt.

\subsection{Testfunksjoner}
Funksjoner i \(V\) som brukes til å formulere den svake formen.
Testfunksjoner lar oss \enquote{teste} en mulig løsning ved å integrere mot disse funksjonene.

\subsection{Basisfunksjoner}
Lokale funksjoner som spenner ut løsningsrommet.
\begin{itemize}
	\item Har kompakt støtte (er null utenfor et lite område).
	\item Oppfyller \(\phi_i(x_j) = \delta_{ij}\) (Kronecker delta)
\end{itemize}
Basisfunksjoner brukes til å konstruere numeriske løsninger. I endelig element-metoden er disse typisk \enquote{hatt-funksjoner} eller polynomer som er positive i små områder og null ellers.

\subsection{Energifunksjon}
La \(F: V \mapsto \mathbb{R}\) være en funksjon som tar inn en funksjon \(v\) og gir ut et reelt tall. Energien til en funksjon \(v\) er gitt ved:
\begin{equation}\label{eq:energy_functional}
	F(v) = \frac{1}{2} \langle v, v \rangle - \langle f, v \rangle
\end{equation}
I mange fysiske problemer representerer dette faktisk den fysiske energien i systemet. Minimering av denne energifunksjonen gir oss løsningen på det tilsvarende differensialligningsmessige problemet, noe som er bakgrunnen for variasjonelle formuleringer.

\subsection{Svak derivasjon}
Svak derivasjon er et konsept som lar oss definere derivasjon av funksjoner som ikke nødvendigvis er klassisk deriverbare.
\begin{definition}{Svak derivasjon}{weak_derivative}
	En funksjon $u \in L^2(\Omega)$ har en svak derivasjon $Du \in L^2(\Omega)$ hvis for alle testfunksjoner $\varphi \in C_c^\infty(\Omega)$ gjelder:
	\[
		\int_\Omega u \frac{\partial \varphi}{\partial x} \, dx = -\int_\Omega Du \varphi \, dx.
	\]
\end{definition}

\subsection{Bilinær form}
\begin{definition}{Bilinær form}{bilinear_form}
	En bilinær form $a: V \times V \to \mathbb{R}$ er en funksjon som er lineær i begge argumentene, dvs. for alle $u,v,w \in V$ og $\alpha, \beta \in \mathbb{R}$ gjelder:
	\begin{align*}
		a(\alpha u + \beta v, w) & = \alpha a(u,w) + \beta a(v,w)  \\
		a(u, \alpha v + \beta w) & = \alpha a(u,v) + \beta a(u,w).
	\end{align*}
\end{definition}

For eksempel i Poisson-ligningen blir den bilinære formen:
\begin{equation}
	a(u,v) = \int_\Omega \nabla u \cdot \nabla v \, dx
\end{equation}

Bilinære former brukes i variasjonelle formuleringer og representerer typisk energien eller arbeidet i det fysiske systemet.

\end{document}